\documentclass[pdflatex,sn-basic,Numbered]{sn-jnl}
\usepackage{graphicx}
\usepackage{amsmath,amssymb,amsfonts}
\usepackage{booktabs}
\usepackage{tabularx}
\begin{document}

\title[Supplementary Information]{Supplementary Information: Fixed-Budget Local-Patch Risk Routing for Industrial Anomaly Detection}
\subtitle{Submitted to Signal, Image and Video Processing}
\author*[1]{\fnm{Jin} \sur{Huang}}\email{614938561@qq.com; \href{https://orcid.org/0009-0005-7489-2018}{ORCID: 0009-0005-7489-2018}}
\author[2]{\fnm{Qisen} \sur{Gao}}\email{1350728839@qq.com}
\affil*[1]{\orgname{Guangxi Vocational Normal University}, \orgaddress{\city{Nanning}, \postcode{530007}, \state{Guangxi Zhuang Autonomous Region}, \country{China}}}
\affil[2]{\orgname{Gansu University of Political Science and Law}, \orgaddress{\city{Lanzhou}, \postcode{730070}, \state{Gansu}, \country{China}}}

\maketitle

\section{Purpose and scope}
This supplement records the unit definitions, paired resampling protocol, controlled baselines, failure cases, online-controller diagnostics, and latency measurement details supporting the main manuscript. All numerical values are from the strict matched-random evaluation unless explicitly labelled as an online or diagnostic result.

\section{Evaluation units and uncertainty}
An evaluation unit is one dataset category--random seed pair. The primary matched comparisons use 45 units for MVTec AD (15 categories $\times$ 3 seeds), 18 units for MPDD (6 $\times$ 3), and 36 units for VisA (12 $\times$ 3). For every unit, the local-routing score and the matched-random-selection score are computed with the same feature extractor, candidate patches, budget, seed, and downstream scoring rule. Reported effects are local minus matched-random values. Confidence intervals are percentile intervals from 10,000 paired bootstrap resamples of evaluation units. The sign-count column reports the number of units with a positive effect.

\begin{table}[ht]
\caption{Primary paired effects (local-patch risk minus matched-random), canonical 25\% strong-routing result}
\label{tab:primary}
\begin{tabularx}{\textwidth}{lrrrrrr}
\toprule
Dataset & $n$ & AUROC $\Delta$ & 95\% CI & $+$ units & Recall $\Delta$ & 95\% CI \\
\midrule
MVTec AD & 45 & 0.1249 & [0.0837, 0.1659] & 33/45 & 0.3877 & [0.2916, 0.4792] \\
MPDD & 18 & 0.1166 & [0.0715, 0.1612] & 16/18 & 0.1975 & [0.1097, 0.2939] \\
VisA & 36 & 0.0128 & [0.0081, 0.0171] & 28/36 & 0.2300 & [0.1336, 0.3214] \\
\bottomrule
\end{tabularx}
\end{table}

The same 99 category--seed rows give a combined AUROC difference of +0.0826 (95\% CI [0.0598, 0.1059]), a Recall@5\%FPR difference of +0.2958 ([0.2347, 0.3551]), and 77/99 units with a positive AUROC effect. Against the additional routing baselines, local-patch risk exceeds the nearest-patch-distance dispersion controller by +0.0704 AUROC ([0.0446, 0.0990]) but only matches the fast-score controller (+0.0017, 95\% CI [-0.0184, 0.0233]); on MPDD and VisA the fast-score controller is stronger.

\section{Budget sensitivity}
The MVTec budget sensitivity uses the same three-seed strict-quota protocol at 10\%, 25\% and 50\% fallback budgets, with 45 category--seed units per budget. Table~\ref{tab:budget} reports the paired AUROC and Recall@5\%FPR differences against the matched-random controller. The AUROC allocation benefit increases with budget, while the Recall benefit decreases; both remain strictly positive. These values describe allocation quality relative to matched-random routing and should not be interpreted as an end-to-end speedup.

\begin{table}[ht]
\caption{MVTec budget sensitivity (local-patch risk minus matched-random) at 10\%, 25\% and 50\% fallback budgets}
\label{tab:budget}
\begin{tabularx}{\textwidth}{lrrrrrX}
\toprule
Budget & $n$ & AUROC $\Delta$ & 95\% CI & Recall $\Delta$ & 95\% CI & $+$ units (AUROC / Recall) \\
\midrule
10\% & 45 & 0.1145 & [0.0821, 0.1484] & 0.4981 & [0.3940, 0.5977] & 39/45, 42/45 \\
25\% & 45 & 0.1249 & [0.0837, 0.1659] & 0.3877 & [0.2916, 0.4792] & 33/45, 42/45 \\
50\% & 45 & 0.1589 & [0.1226, 0.1954] & 0.2479 & [0.1834, 0.3110] & 39/45, 35/45 \\
\bottomrule
\end{tabularx}
\end{table}

\section{Baseline and absolute-score context}
On the 45 MVTec units, the full detector obtains AUROC 0.9392. At the primary 25\% local budget, the strict local route obtains AUROC 0.8343 and the matched-random route 0.7094. These absolute values are not presented as superiority over the full detector; they establish that the contribution is allocation quality at a fixed reduced local budget. A PaDiM-style baseline obtains AUROC 0.8982 and Recall@5\%FPR 0.6544. A controlled PatchCore-like patch-level farthest-point coreset detector obtains AUROC 0.9433 and Recall@5\%FPR 0.7847. The latter is an implementation-controlled approximation, not a claim of exact parity with the official PatchCore release.

\section{Failure cases and boundary conditions}
Three checks were retained as negative evidence. First, a multi-scale rank-fusion variant underperformed matched-random routing in a three-category smoke test, indicating that additional ranking signals can miscalibrate the budget allocation. Second, a lightweight cost filter underperformed the stronger patch-level coreset baseline. Third, the conformal-threshold diagnostic produced a realized fallback rate of 54.9\% for a nominal 25\% target and is therefore not used as a budget guarantee. These tests motivate the conservative claim that routing improves selection quality under a fixed experimental budget, not that every controller or every online policy is robust.

\section{Online-prefix sensitivity}
The online prefix controller was evaluated separately because it introduces a temporal prefix and fallback policy. On the 45 MVTec units, realized fallback was 25.34\%, with AUROC effect 0.0331 (95\% CI [0.0045, 0.0601]) and Recall effect 0.0804 (95\% CI [0.0211, 0.1346]). The same controller produced AUROC/Recall effects of 0.0890/0.0643 on MPDD seed 17 and 0.00004/$-0.0408$ on VisA seed 17. With a VisA buffer of 128, the effects were 0.00124 (95\% CI [$-0.00113$, 0.00339]) and 0.0750 (95\% CI [0.0189, 0.1317]). These results are why online generalization is reported as configuration-sensitive rather than as a universal improvement.

\section{Latency measurement}
Latency was measured on the same machine with warm-up excluded from the reported percentile. Across 15 MVTec categories and 6,000 images with batch size one and CUDA synchronization, the Fast path had mean/P95 2.3292/2.9298 ms, the Full path 3.6269/4.0860 ms, and the Risk path (probe plus full) 5.1899/6.7302 ms. The Risk path includes both local and full operations and is a boundary diagnostic; it is not an end-to-end deployment benchmark. Batch throughput checks were limited to three categories and are treated as an extension rather than a primary claim.

\section{Reproducibility record}
The submission package records the random seeds, category lists, budget values, derived tables and figure source data. A submission-matched code release, including the environment specification and regeneration scripts, will be archived at \url{https://github.com/cdhuangjin/fixed-budget-patch-routing}. No private path is presented as a reproducibility link.

\end{document}
